\section{Related Work}
\label{sec:related}

\paragraph{Looped and recurrent-depth language models.}
Weight sharing across depth goes back to Universal Transformers \citep{dehghani2019universal}.
Recent work scales it to pretrained language models: a recurrent-depth model trained with a
random number of iterations can spend more loops at test time \citep{geiping2025recurrent};
Ouro trains a looped model with an entropy-regularised exit distribution over four loops
\citep{zhu2025ouro}; Dory-1.0 loops a middle block between fixed input and output blocks and
supports a last-loop KV cache \citep{nvidia2026dory}; Nanbeige4.2 runs its whole stack twice and
reports little benefit from more passes \citep{nanbeige2026}. Theory and controlled studies relate
loops to multi-step reasoning \citep{saunshi2025latent,kohli2026loop}, and fine-tuned looped
models improve compositional tool calling \citep{popescu2026toolcalling}. These works fix the
loop count or use the model's own halting rule at inference. Our work differs from
\citet{popescu2026toolcalling}, the closest to it, in question and in protocol: we study where
to spend loops rather than whether they help; calls are executed one turn at a time rather than
generated as one sequence; every allocation policy is compared with random mixes of fixed depths
and with a sampling control at equal compute, under batched serving; and confirmatory results
are fixed in advance and measured on held-out requests.

\paragraph{Adaptive depth: halting, early exit and routing.}
Adaptive computation time \citep{graves2016adaptive} and PonderNet \citep{banino2021pondernet}
learn when a recurrent computation should stop; Ouro's exit gate follows this line
\citep{zhu2025ouro}. Diagnoses of learned gates in looped transformers find that they are not
Pareto-optimal and that their depth use does not track problem difficulty \citep{popescu2026halting,olszowiec2026monitoring}, and Ouro's per-loop KV caches remove most of the
savings of exiting early \citep{mishra2026dynamic}. In standard
transformers, early exit stops at an intermediate layer when the model is confident
\citep{schuster2022calm,elhoushi2024layerskip}, and learned routers skip or repeat layers per
token or per input
\citep{raposo2024mod,he2025routertuning,luo2025flexidepth,heakl2026drllm,li2025cola}, in one
case choosing exit layers with reinforcement learning \citep{zhu2026lede}.
Mixture-of-Recursions trains a per-token router over recursion depth during pretraining
\citep{bae2025mor}. Our router differs in three ways: it acts once per request rather than per
token, so batched serving keeps its savings; it is fitted after training on a frozen looped model;
and it is evaluated on multi-step tool use with executed calls.

\paragraph{Speculative decoding with loops.}
Because a shallow loop's prediction is a prefix of the full computation, looped models can
draft tokens at few loops and verify them at full depth without changing the output
\citep{leviathan2023speculative,geiping2025recurrent}. LoopSpec pipelines this and reports about
3$\times$ faster batch-size-1 decoding on Ouro-2.6B \citep{cho2026loopspec}, with related
schemes in \citet{ha2026wavefront}, \citet{li2026das} and \citet{liu2026loopahead}. These methods
reduce latency at a fixed depth and leave the output unchanged; they do not decide the depth.
They are complementary to ours, and their observation that one-loop drafts usually match the
four-loop output agrees with our finding that loops disagree mainly on decision tokens.
Cross-loop distillation instead trains shallow loops to imitate the deepest one
\citep{wang2026loopopd}.

\paragraph{Allocating test-time compute.}
Spending more inference compute on harder inputs can beat a larger model at equal cost
\citep{snell2024scaling}, and learned predictors can allocate decoding budget per input
\citep{damani2024learning}. Width-based methods sample several answers and vote
\citep{wang2023selfconsistency}. Model routing and cascades choose between models of different
sizes from the input or from a cheap model's confidence
\citep{chen2024frugalgpt,ong2025routellm,gupta2024cascades}. We allocate depth within one looped
model, compare depth with width at equal compute, and show that the input alone does not reveal
which requests need more depth.

\paragraph{Tool use.}
Language models learn to call tools from demonstrations and synthetic data
\citep{schick2023toolformer,patil2024gorilla,liu2024apigen} and are evaluated on single and
parallel calls \citep{patil2025bfcl} and on nested sequences whose calls depend on earlier
results \citep{basu2024nestful}. NESTful's math portion is built from the MathQA test split
\citep{amini2019mathqa}, whose problems derive from AQuA-RAT \citep{ling2017aqua}; we use the
same construction on MathQA's training split for fine-tuning and audit its overlap with
validation and test.
